% Generated draft: requires mathematical and source review.
\documentclass[11pt]{article}
\usepackage[margin=25mm]{geometry}
\usepackage{fontspec}
\setmainfont{Latin Modern Roman}
\usepackage{amsmath,amssymb,mathtools,unicode-math}
\setmathfont{Latin Modern Math}
\usepackage{enumitem,needspace,xurl,hyperref}
\setlength{\parindent}{0pt}
\setlength{\parskip}{5pt}
\setlength{\emergencystretch}{3em}
\newcommand{\N}{\mathbb{N}}
\newcommand{\Z}{\mathbb{Z}}
\newcommand{\Q}{\mathbb{Q}}
\newcommand{\R}{\mathbb{R}}
\newcommand{\C}{\mathbb{C}}
\newcommand{\F}{\mathbb{F}}
\newcommand{\A}{\mathbb{A}}
\newcommand{\PP}{\mathbb{P}}
\newcommand{\E}{\mathbb{E}}
\newcommand{\eps}{\varepsilon}
\newcommand{\dd}{\,\mathrm{d}}
\begin{document}
\section*{UnsolvedMath ID: 1101218}
\textbf{Source problem code:} catalogue entry
\subsection{Definitions and mathematical statement}
Let $S_g$ be a closed connected orientable surface of genus $g\geq3$, and write $\Gamma=\operatorname{MCG}(S_g)$. The source does not specify whether this denotes orientation-preserving mapping classes or the extended mapping class group allowing orientation reversal. Fix the intended convention and a word metric $d$ from a finite generating set of $\Gamma$.

A self quasi-isometry is a map $f:\Gamma\to\Gamma$ for which there exist $A\geq1$ and $B\geq0$ such that
\[
A^{-1}d(x,y)-B\leq d(f(x),f(y))\leq A d(x,y)+B
\]
for all $x,y\in\Gamma$, and every point of $\Gamma$ is within $B$ of $f(\Gamma)$. Two such maps are equivalent when $\sup_{x\in\Gamma}d(f(x),h(x))<\infty$. The group of equivalence classes under composition is $\operatorname{QI}(\Gamma)$, independent of the chosen finite generating set. Left translation $L_a(x)=ax$ defines the natural homomorphism
\[
\iota_g:\Gamma\longrightarrow\operatorname{QI}(\Gamma),\qquad a\longmapsto[L_a].
\]
\textbf{Question.} For every integer $g\geq3$, is $\iota_g$ an isomorphism? Equivalently, does every self quasi-isometry of $\Gamma$ lie at finite uniform distance from left translation by a unique element of $\Gamma$?

The supplied summary reports an affirmative answer as the closed-surface case of Behrstock--Kleiner--Minsky--Mosher, Corollary 10.1. Its application depends on the mapping-class convention: orientation-reversing mapping classes also induce self quasi-isometries of the orientation-preserving mapping class group, so the distinction between the full and orientation-preserving groups is substantive and is not resolved by the source notation alone.
\subsection{Short English statement}
Is every quasi-isometry of the mapping class group of a closed surface of genus at least three a bounded distance from a unique left translation? The supplied affirmative theorem citation requires fixing whether the mapping class group includes orientation reversal.
\subsection{Sources}
\begin{itemize}
\item M. Bestvina, \emph{Questions in Geometric Group Theory} (2004), Question 12.18, PDF page 23: \url{https://www.math.utah.edu/~bestvina/eprints/questions-updated.pdf}; manual transcription, accessed 2026-07-29.
\item Behrstock--Kleiner--Minsky--Mosher, \emph{Geometry \& Topology} 16 (2012), 781--888, Corollary 10.1, DOI 10.2140/gt.2012.16.781, \url{https://doi.org/10.2140/gt.2012.16.781}. This is the supplied research-summary citation and claimed closed-surface application.
\item Local registry \texttt{data/unsolvedmath/problems.json}, permanent ID $1101218$, source code catalogue entry, \emph{Questions in Geometric Group Theory --- Q 12.18}. The record labels this solved and SOLVED-IN-LITERATURE; its dated source-list background retains NEEDS\_REVIEW status and rights qualifications.
\end{itemize}
\end{document}
